%%
%% MolTrustBench manuscript using the Oxford University Press authoring
%% template. Author metadata remains placeholder text until submission.
%%

\documentclass[unnumsec,webpdf,contemporary,large,numbered]{oup-authoring-template}

\usepackage[T1]{fontenc}
\usepackage{lmodern}
\usepackage{booktabs}
\usepackage{float}
\usepackage{graphicx}
\usepackage{url}
\urlstyle{same}
\Urlmuskip=0mu plus 1mu
\makeatletter
\let\@unusedoptionlist\@empty
\@numbibtrue
\NAT@numberstrue
\makeatother
\setcitestyle{numbers,square,comma}

\graphicspath{{figures/}{../../results/figures/}}

\newcommand{\artifact}[1]{\begingroup\Urlmuskip=0mu plus 1mu\path{#1}\endgroup}
\newcommand{\placeholderfigure}[1]{%
  \begin{center}
  \fbox{\begin{minipage}[c][1.65in][c]{0.88\linewidth}
  \centering #1
  \end{minipage}}
  \end{center}
}
\newcommand{\maybefigure}[2][]{%
  \IfFileExists{figures/#2}{%
    \includegraphics[#1]{figures/#2}%
  }{%
    \IfFileExists{../../results/figures/#2}{%
      \includegraphics[#1]{../../results/figures/#2}%
    }{%
      \placeholderfigure{Planned figure artifact: \artifact{#2}}%
    }%
  }%
}

\begin{document}

\journaltitle{Briefings in Bioinformatics}
\DOI{DOI added during production}
\copyrightyear{2026}
\pubyear{2026}
\vol{XX}
\issue{x}
\access{Published: Date added during production}
\appnotes{Problem Solving Protocol}

\firstpage{1}

\title[MolTrustBench]{Benchmarking without Time Travel: Public-Exposure Auditing and Temporal Generalization in Molecular AI}

\author[1,$\ast$]{SUBMISSION\_METADATA\_REQUIRED\_AUTHOR\_1}
\author[1]{SUBMISSION\_METADATA\_REQUIRED\_AUTHOR\_2}
\author[1]{MolTrustBench Consortium}

\address[1]{\orgdiv{SUBMISSION\_METADATA\_REQUIRED\_DEPARTMENT}, \orgname{SUBMISSION\_METADATA\_REQUIRED\_INSTITUTION}, \orgaddress{\street{SUBMISSION\_METADATA\_REQUIRED\_STREET}, \postcode{SUBMISSION\_METADATA\_REQUIRED\_POSTCODE}, \state{SUBMISSION\_METADATA\_REQUIRED\_STATE}, \country{SUBMISSION\_METADATA\_REQUIRED\_COUNTRY}}}

\corresp[$\ast$]{Corresponding author. \url{mailto:submission-metadata-required@example.invalid}}

\received{Date}{0}{2026}
\revised{Date}{0}{2026}
\accepted{Date}{0}{2026}

\abstract{Molecular AI benchmarks are commonly interpreted as prospective generalization tests, but held-out molecules, scaffolds, close chemical neighbors, and assay-derived labels can already be publicly observable before a relevant benchmark cutoff. This creates a temporal-validity problem that is distinct from conventional ADMET, out-of-distribution, class-imbalance, or activity-cliff benchmarking. We introduce MolTrustBench, a reproducible audit protocol and reporting standard that builds ChEMBL release-time indexes, annotates exact, scaffold, and nearest-neighbor public exposure, constructs exposure-aware evaluation slices, audits assay provenance, and emits benchmark trust cards. In a first ChEMBL24/27/30/33/36 audit of MoleculeNet BBBP, ClinTox, BACE, and Tox21, public exposure is measurable at multiple structural levels and varies strongly by task. Exposure-aware evaluation shows task- and model-family-dependent score sensitivity, while density, temporal, and train-label-shuffle controls define interpretation boundaries. ChEMBL36 assay provenance reveals duplicate records, unit heterogeneity, threshold sensitivity, and provenance-level label discordance for benchmark-overlapping and endpoint-native toxicology tasks. MolTrustBench reports observable public-exposure lower bounds, not model-specific training inclusion or memorization, and provides machine-readable supplement tables, data-freeze manifests, and trust-card templates for benchmark reporting.}

\keywords{molecular AI, benchmark trustworthiness, public exposure, temporal validity, ChEMBL, assay provenance}

\keywords[Abbreviations]{ADMET: absorption, distribution, metabolism, excretion, and toxicity; NN: nearest neighbor; TDC: Therapeutics Data Commons; AUROC: area under the receiver operating characteristic curve}

\boxedtext{Key Points}{
\begin{itemize}
\item MolTrustBench audits benchmark temporal validity rather than proposing another ADMET leaderboard.
\item ChEMBL release-time indexes provide observable lower bounds for molecule, scaffold, and nearest-neighbor public exposure.
\item Exposure-adjusted evaluation and assay-provenance diagnostics reveal benchmark trustworthiness limitations that are invisible in single leaderboard scores.
\end{itemize}}

\maketitle

\section{Introduction}

Molecular AI benchmarks increasingly serve as proxies for prospective drug-discovery and toxicology generalization. This interpretation assumes that held-out molecules represent chemistry that was not publicly observable at the relevant evaluation boundary. In the foundation-model era, however, a molecule can be absent from a supervised training split while still being publicly observable through ChEMBL, PubChem, ZINC, benchmark repositories, or assay-derived records before evaluation \cite{chembl2012,chembl2019,chembl2023,chemblDownloads2025,pubchem2019,zinc2020}. This distinction matters because a leaderboard score can mix prospective generalization with interpolation over publicly visible chemistry.

Existing molecular benchmark work has made substantial progress on data scarcity, out-of-distribution generalization, scaffold splits, class imbalance, activity cliffs, and ADMET reliability \cite{moleculenet2018,tdc2021,ogb2020,wallach2018,sheridan2013,moleculeace2023,deeptox2016,boom2025}. Recent molecular representation and foundation-model work has also treated benchmark overlap with public molecular corpora as a concrete evaluation-control issue \cite{yang2019representation,hu2020pretraining,mole2024}. MolTrustBench addresses a different validity question: whether the benchmark itself should be trusted as a temporal evaluation of chemistry not observable in an indexed public release history before a cutoff. We therefore use conservative terminology throughout. A ChEMBL release date is an observability date, not a model-specific pretraining date; public exposure is an observable lower bound, not proof that a particular model used or recalled a molecule.

\paragraph{Interpretation limits.} Public observability does not establish model training inclusion, memorization, or inappropriate benchmark use. MolTrustBench therefore audits an observable public-exposure lower bound and reports exposure-adjusted score sensitivity. These outputs should be interpreted as benchmark-trustworthiness diagnostics, not as accusations about a particular model or checkpoint.

MolTrustBench contributes a reproducible audit framework for benchmark trustworthiness. The framework builds historical release indexes, annotates exact molecule, scaffold, and nearest-neighbor exposure, constructs exposure-aware and temporal-future evaluation slices, audits assay provenance from ChEMBL SQLite, and generates benchmark trust cards. This first MolTrustBench audit uses real ChEMBL release indexes, four MoleculeNet case studies, endpoint-native toxicology pilots, negative controls, and machine-readable trust-card outputs to illustrate the protocol rather than to introduce a new ADMET leaderboard.

\begin{table*}[!t]
\centering
\caption{\textbf{Novelty boundary for MolTrustBench.} The paper is framed as a benchmark-trustworthiness protocol and reporting standard, not as a replacement for existing molecular performance benchmarks.}
\label{tab:novelty-boundary}
\small
\begin{tabular*}{\textwidth}{@{\extracolsep{\fill}}p{0.30\textwidth}p{0.62\textwidth}}
\toprule
Adjacent benchmark area & MolTrustBench boundary \\
\midrule
ADMET reliability benchmarks & Audits temporal observability and assay provenance rather than ranking ADMET predictors under scarcity, imbalance, or endpoint difficulty. \\
Molecular OOD and scaffold-split benchmarks & Reports whether test chemistry was publicly observable before a cutoff; this is complementary to distribution-shift or scaffold generalization tests. \\
Activity-cliff benchmarks & Treats scaffold and nearest-neighbor exposure as public-observability views, not as a substitute for local structure--activity difficulty. \\
Contamination or memorization studies & Provides observable exposure lower bounds and exposure-adjusted evaluations; it does not prove training-set inclusion or memorization. \\
Foundation-model evaluations & Uses frozen embedding wrappers only as protocol applicability checks, not as evidence that a checkpoint used ChEMBL. \\
Drug-discovery agent benchmarks & Does not evaluate autonomous discovery agents; it audits whether benchmark evidence is temporally and provenance-aware. \\
\bottomrule
\end{tabular*}
\end{table*}

\begin{figure*}[!t]
\centering
\maybefigure[width=0.92\textwidth,keepaspectratio]{workflow_schematic.pdf}
\caption{\textbf{MolTrustBench workflow.} MolTrustBench converts benchmark molecules and labels into auditable public-observability, exposure-adjusted evaluation, assay-provenance, and trust-card artifacts. Benchmark molecules are standardized, linked to ChEMBL release-time compound and scaffold indexes, annotated for exact, scaffold, and nearest-neighbor public exposure, evaluated under exposure-aware splits, and summarized in benchmark trust cards. ChEMBL release dates are treated as public observability dates, not as evidence that a specific model used ChEMBL during pretraining.}
\label{fig:workflow}
\end{figure*}

\section{Methods}

\subsection{Release-time public-exposure index}

We use historical ChEMBL releases as public-observability anchors. For each release, MolTrustBench records standard InChIKey, standardized SMILES, release metadata, and Bemis--Murcko scaffold information \cite{chemblDownloads2025,bemis1996}. Compound-level and scaffold-level indexes are collapsed to first-seen release and first-seen date. The initial index covers CHEMBL24, CHEMBL27, CHEMBL30, CHEMBL33, and CHEMBL36.

\subsection{Benchmark standardization and exposure annotation}

Benchmark molecules are standardized with RDKit when available \cite{rdkit}. Each molecule is annotated with exact public exposure, scaffold public exposure, nearest-neighbor exposure at Tanimoto thresholds 0.6, 0.7, 0.8, and 0.9, maximum nearest-neighbor similarity before the cutoff, earliest exact release, earliest scaffold release, exposure category, and an exposure audit score. The main cutoff in this audit is CHEMBL30, corresponding to February 2022 in the ChEMBL release table.

\subsection{Exposure-adjusted evaluation}

For the current performance evidence, MolTrustBench compares standard splits against exact-removed, scaffold-removed, NN-removed, density-matched exposure-removed, and observable temporal-future subsets. Classical Morgan-fingerprint baselines use logistic regression and an XGBoost-style gradient-boosted classifier fallback \cite{ecfp2010,xgboost2016,scikitlearn2011}, while sequence-neural baselines include SMILES-CNN, SMILES-GRU, and SMILES-Transformer families \cite{weininger1988,honda2019smilestransformer}. Regression runs report RMSE, MAE, and $R^2$ under the same exposure-aware slice logic. These models are used to test score sensitivity under exposure-aware evaluation; they are not treated as evidence of foundation-model training exposure.

Train-label-shuffle controls permute training labels before fitting sequence baselines, then compare null-control exposure-adjusted deltas with observed-label deltas. This analysis is used as a negative control for score sensitivity, not as evidence that public exposure caused a specific model behavior.

For endpoint-native temporal evaluation, MolTrustBench also constructs ChEMBL-derived hERG, CYP, and DILI/hepatotoxicity tasks from release-indexed bioactivity or text-evidence records. These tasks use the same exposure-aware split logic but are interpreted as endpoint-specific temporal-validity and provenance checks. Strict exposure-removed slices that become too small or one-class are retained as expected limitations rather than promoted to performance results. Frozen foundation-model embeddings are treated as checkpoint-light wrapper case studies, not evidence that a specific checkpoint used ChEMBL during pretraining.

\subsection{Assay provenance}

Assay-provenance auditing uses ChEMBL36 SQLite activity tables to recover assay identifiers, document identifiers, targets, assay types, standard measurement fields, units, relations, pChEMBL values when present, confidence scores, and source release. MolTrustBench summarizes duplicate compound counts, discordant threshold-derived label calls, unit inconsistency, threshold sensitivity, document-level heterogeneity, and a derived-label discordance score. The current evidence includes MoleculeNet-overlapping BBBP and ClinTox molecules plus endpoint-native ChEMBL-derived hERG, CYP1A2, CYP2C19, and DILI/hepatotoxicity endpoints.

\section{Results}

\subsection{MolTrustBench establishes a release-time index for public-exposure auditing}

MolTrustBench first constructs a release-time index that converts public chemical database history into an auditable benchmark timeline. We indexed five ChEMBL releases spanning CHEMBL24, CHEMBL27, CHEMBL30, CHEMBL33, and CHEMBL36, with release dates from May 2018 to July 2025 and standardized molecule counts ranging from 1.82 million to 2.85 million records (Table~\ref{tab:chembl-releases}). For each release, the pipeline records standard InChIKey, standardized SMILES, release metadata, and Bemis--Murcko scaffold information, then collapses records into earliest observable release indexes at the compound and scaffold levels. This design intentionally treats ChEMBL release date as an observability date, not as proof that any specific model used ChEMBL during pretraining.

\begin{table}[!t]
\caption{\textbf{ChEMBL release index used for the first MolTrustBench audit.}}
\label{tab:chembl-releases}
\begin{tabular*}{\columnwidth}{@{\extracolsep\fill}llllr}
\toprule
Release & Date & DOI suffix & Input & Std. molecules \\
\midrule
CHEMBL24 & 2018-05 & database.24 & chemreps & 1,820,035 \\
CHEMBL27 & 2020-05 & database.27 & chemreps & 1,941,410 \\
CHEMBL30 & 2022-02 & database.30 & chemreps & 2,136,187 \\
CHEMBL33 & 2023-05 & database.33 & chemreps & 2,372,674 \\
CHEMBL36 & 2025-07 & database.36 & chemreps & 2,854,815 \\
\bottomrule
\end{tabular*}
\end{table}

\subsection{Public exposure is measurable across benchmark tasks}

The real-data audit finds measurable public exposure across all four completed MoleculeNet tasks, supporting the claim that benchmark trustworthiness should be audited before interpreting held-out molecules as prospectively unseen chemistry. At the CHEMBL30 cutoff, exact exposure was 0.838 for BBBP, 0.474 for ClinTox, 0.297 for BACE, and 0.921 for Tox21 (Table~\ref{tab:coverage}). The four-task audit covers 12,290 benchmark molecules, with ChEMBL-index-unobserved exact counts ranging from 331 in BBBP to 1,063 in BACE. These values are not uniform across tasks: BACE contains a larger exact-unobserved subset, whereas Tox21 is already highly exact-exposed under the same cutoff. This heterogeneity is the point of the audit: a benchmark trust card should report the public-exposure profile of each task rather than assuming that all benchmark test sets carry the same temporal meaning.

\begin{figure}[!t]
\centering
\maybefigure[width=0.94\columnwidth,keepaspectratio]{exposure_heatmap_coverage.pdf}
\caption{\textbf{Multi-level public-exposure landscape.} Exact, scaffold, and nearest-neighbor exposure rates vary across MoleculeNet BBBP, ClinTox, BACE, and Tox21 at the CHEMBL30 cutoff. The heatmap shows that exact molecule occurrence is only one exposure view: scaffold and NN exposure can reveal additional observable exposure lower bounds, while some tasks such as Tox21 are already highly exact-exposed. The result motivates per-benchmark trust cards instead of a single universal exposure label. Denominators are the audited benchmark molecule counts in Table~\ref{tab:coverage}.}
\label{fig:coverage-heatmap}
\end{figure}

\begin{table}[!t]
\caption{\textbf{Audit-only exposure coverage at the CHEMBL30 cutoff.}}
\label{tab:coverage}
\begin{tabular*}{\columnwidth}{@{\extracolsep\fill}lrrrr}
\toprule
Task & Molecules & Exact & Scaffold & NN@0.8 \\
\midrule
BBBP & 2,039 & 0.838 & 0.878 & 0.889 \\
ClinTox & 1,480 & 0.474 & 0.737 & 0.599 \\
BACE & 1,513 & 0.297 & 0.427 & 0.394 \\
Tox21 & 7,258 & 0.921 & 0.753 & 0.959 \\
\bottomrule
\end{tabular*}
\end{table}

\subsection{Scaffold and nearest-neighbor exposure extend beyond exact matching}

Exact molecule matching is only the most conservative exposure test, and the scaffold and nearest-neighbor views provide additional public-observability views. For BBBP, exact exposure at CHEMBL30 was 0.838, while scaffold exposure was 0.878 and nearest-neighbor exposure at Tanimoto $\geq 0.8$ was 0.889. ClinTox showed a larger gap between exact and structural exposure: exact exposure was 0.474, but scaffold exposure reached 0.737 and NN@0.8 reached 0.599. BACE showed the same ordering at lower absolute rates, with exact exposure of 0.297, scaffold exposure of 0.427, and NN@0.8 exposure of 0.394. Tox21 is an important counterexample rather than an inconvenience: exact exposure was already very high at 0.921, while scaffold exposure was 0.753 and NN@0.8 was 0.959. These results indicate that public-exposure auditing is multi-level and task-specific, not a single monotonic score.

\subsection{Exposure-adjusted evaluation can change benchmark interpretation}

Exposure-adjusted evaluation shows that standard split scores and exposure-removed subset scores can diverge enough to change how benchmark performance is interpreted. On BBBP, Morgan-logistic regression achieved AUROC 0.879 on the standard split and 0.946 on the exact-removed split; Morgan-XGB similarly changed from 0.885 to 0.976. ClinTox moved in the opposite direction: Morgan-logistic regression changed from AUROC 0.836 on the standard split to 0.565 on the exact-removed split, and Morgan-XGB changed from 0.784 to 0.370. These results are not interpreted as a universal performance drop. Instead, exposure-aware subsets can produce different estimates of model performance, with direction and magnitude depending on task composition, label structure, model class, and exposure-removed subset size.

\begin{figure}[!t]
\centering
\maybefigure[width=0.94\columnwidth,keepaspectratio]{exposure_delta_ci.pdf}
\caption{\textbf{Exposure-adjusted score deltas.} The main heatmap shows rows with at least five matched runs and non-degenerate seed-replicate uncertainty. Deltas are computed as standard-score minus comparison-score for matched standard and exposure-aware slices; open markers indicate intervals excluding zero only when repeated-run uncertainty is available. Single-run or deterministic-repeat rows are retained in the source tables as descriptive estimates without interval markers. The figure is evidence of benchmark score sensitivity under public-exposure auditing; it does not prove model-specific training exposure or label recall.}
\label{fig:performance-drop}
\end{figure}

\subsection{Density-matched and temporal-future controls define the current evidence boundary}

The first controls address two interpretation caveats: public exposure may track chemical-space density, and release dates are not model-specific training dates. Density-matched exposure-removed controls did not eliminate score differences in the evaluated tasks, but exposure-removed subsets were often much smaller than standard test sets. Temporal-future splits remain a narrower control: BBBP had 34 temporal-future test molecules, ClinTox had only 8, and BACE/Tox21 temporal-future subsets were too small or single-class for stable classification metrics. Metric-computable rows with very low minority-class counts are treated as boundary evidence rather than stable performance estimates. These controls support the framework logic, while temporal-generalization evidence remains a current boundary rather than a final deployment claim.

\subsection{Endpoint-native ChEMBL tasks add temporal-validity checks}

Endpoint-native ChEMBL temporal tasks add a stress test because labels and release history are derived from ChEMBL bioactivity records rather than mapped retrospectively from MoleculeNet labels. hERG and CYP sequence baselines use the same standard, exposure-removed, density-matched, and temporal-future logic where valid. The summaries show endpoint-specific temporal drops and class-degenerate exposure-removed slices, which the protocol records as benchmark-usability findings rather than hidden failures.

\subsection{Frozen FM embeddings provide supplement-first applicability checks}

MolFormer and ChemBERTa-100M frozen-embedding runs are reported in supplement tables as applicability checks, not as a foundation-model leaderboard \cite{molformer2022,chemberta2020}. They show that the audit contract can wrap staged checkpoints without implying that either checkpoint used ChEMBL during pretraining or generalizing to all molecular foundation models.

\subsection{Train-label-shuffle controls provide a negative-control check}

Train-label-shuffle controls test whether exposure-adjusted deltas remain larger than a null training-label baseline. Across 22 observed-vs-shuffle comparisons, most observed deltas are larger than shuffled-label deltas, while four remain null-sensitive. The supplement reports the full null-control table and expected one-class limitations; the result bounds interpretation rather than serving as causal evidence that public exposure changed any model.

\subsection{Model-family sensitivity shows that exposure effects are not confined to one baseline}

The model-family sweep extends exposure-adjusted evaluation beyond classical fingerprints while keeping the interpretation conservative. Across BBBP, ClinTox, BACE, and Tox21, the artifacts include Morgan, MLP, SMILES-CNN, SMILES-GRU, and SMILES-Transformer results under valid standard, exposure-removed, density-matched, and temporal-future slices. Repeated BACE/Tox21 sequence runs separate stable deltas from uncertain deltas, and ESOL regression adds the same audit logic for RMSE. The goal is not to rank models, but to show that benchmark trustworthiness should be reported by exposure slice and model family rather than by a single leaderboard number.

\begin{figure}[!t]
\centering
\maybefigure[width=0.98\columnwidth,keepaspectratio]{bace_tox21_sequence_model_family.pdf}
\caption{\textbf{Sequence-family exposure sensitivity.} BACE/Tox21 sequence-family baselines show that exposure-adjusted score changes are architecture- and task-dependent. In panel A, heatmap colors encode standard-minus-comparison AUROC deltas; open circles denote intervals excluding zero and x marks intervals crossing zero. In panel B, blue bars denote BACE comparison-slice sizes and orange bars denote Tox21 comparison-slice sizes. The integrated sweep covers SMILES-CNN, SMILES-GRU, and SMILES-Transformer baselines after run-identifier deduplication. These models are sequence baselines for exposure-adjusted evaluation, not foundation-model pretraining evidence.}
\label{fig:model-family}
\end{figure}

\subsection{Assay provenance reveals label-source heterogeneity beyond molecule exposure}

MolTrustBench also audits whether benchmark-label molecules have heterogeneous public assay records. ChEMBL36 SQLite extraction identified activity records for 1,289 BBBP molecules and 654 ClinTox molecules, with duplicate records, unit heterogeneity, threshold sensitivity, and discordant threshold-derived label calls among molecules with public activity records. The endpoint-native hERG, CYP1A2, CYP2C19, and DILI/hepatotoxicity cases extend this provenance check. Together, these results show that benchmark trustworthiness depends not only on molecule exposure, but also on assay source, measurement units, censoring relations, thresholding choices, text-label evidence, and document provenance.

\begin{figure}[!t]
\centering
\maybefigure[width=0.94\columnwidth,keepaspectratio]{assay_conflict_map.pdf}
\caption{\textbf{Assay-provenance heterogeneity map.} ChEMBL36 SQLite assay provenance reveals duplicate activity records, unit heterogeneity, threshold sensitivity, and derived-label discordance for benchmark molecules with public assay records. BBBP and ClinTox rows represent public ChEMBL assay heterogeneity among overlapping molecules, while hERG, CYP1A2, CYP2C19, and DILI/hepatotoxicity provide endpoint-native ChEMBL-derived provenance cases. Denominators are molecules with matched public activity records, and flag definitions are given in Table S26. The figure supports label-source reliability auditing, not a claim that original benchmark labels are simply wrong.}
\label{fig:assay-conflict}
\end{figure}

\subsection{Trust cards convert the audit into a reusable reporting standard}

The final result is a reporting layer that turns exposure annotations, exposure-adjusted metrics, assay provenance, and sanity checks into benchmark trust cards. For each audited task, MolTrustBench records exposure rates, exact-unobserved counts, cutoff release, assay-provenance availability, and audit caveats. The card framing is important because this article is not proposing a new ADMET leaderboard. Instead, it asks benchmark users to report whether test molecules and labels were publicly observable before a relevant cutoff and whether performance estimates are stable under exposure-aware slices.

\begin{figure}[!t]
\centering
\maybefigure[width=0.94\columnwidth,keepaspectratio]{trust_card_examples.pdf}
\caption{\textbf{Benchmark trust-card examples.} Benchmark trust cards summarize public-exposure rates, exposure-removed subset sizes, exposure-adjusted score evidence, assay-provenance availability, and audit caveats for individual benchmark tasks. The card format turns MolTrustBench outputs into a reporting standard so benchmark users can inspect temporal-validity and provenance limitations before interpreting leaderboard scores as prospective generalization. Machine-readable schema and JSON examples are provided in Table S28 and the supplement trust-card directory.}
\label{fig:trust-cards}
\end{figure}

\begin{table}[!t]
\caption{\textbf{Minimum reporting checklist.}}
\label{tab:checklist}
\begin{tabular*}{\columnwidth}{@{\extracolsep\fill}lp{0.55\columnwidth}}
\toprule
Item & Required report field \\
\midrule
Public cutoff & Release or date used as the public-observability boundary \\
Exposure rates & Exact, scaffold, and nearest-neighbor public exposure \\
Exposure-removed subset & Exact-removed, scaffold-removed, and NN-removed test sizes \\
Assay provenance & Assay IDs, document IDs, units, relations, and derived-label discordance scores when available \\
Model metadata & Claimed pretraining corpus and cutoff, or explicit unknown status \\
Adjusted scores & Standard, exposure-removed, exposed, temporal-future, and density-matched scores \\
Limitations & Whether evidence is audit-only, performance-tested, or endpoint-native \\
\bottomrule
\end{tabular*}
\end{table}

\section{Discussion}

MolTrustBench reframes molecular benchmark evaluation as a trustworthiness audit. The central finding is not that any model used a specific molecule or label, but that commonly used benchmark tasks contain measurable public-observability and assay-provenance limitations that should be reported before leaderboard scores are interpreted as prospective generalization. This framing separates three questions that are often conflated: whether chemistry was publicly observable, whether model performance changes under exposure-aware evaluation, and whether assay-derived labels are stable enough to support benchmark claims.

The evidence package is intentionally scoped. The four-task audit supports the molecule, scaffold, and nearest-neighbor public-observability claims. Exposure-adjusted evaluation is supported by the core tasks plus endpoint-native, sequence-family, regression, frozen-embedding, and train-label-shuffle extensions summarized in the supplement. Assay-provenance claims are supported by ChEMBL36 SQLite extraction plus endpoint-native toxicology evidence. The reporting-standard contribution is supported by the workflow schematic, trust cards, sanity checks, data-freeze records, and machine-readable supplement tables S24--S29.

The main limitations are also explicit: SIDER remains deferred; Lipophilicity, CYP2C9, DILI, and some label-shuffle slices become too small or class-degenerate under strict exposure-aware filtering; four observed-vs-null comparisons remain null-sensitive; temporal-future subsets are small for some tasks; frozen FM rows are wrapper applicability evidence; and model-specific pretraining cutoffs remain unknown for many public checkpoints.

These limitations are audit outputs rather than failures to be hidden. Exposure-removed subsets that become too small or class-degenerate under strict filtering are evidence about benchmark usability under exposure-adjusted evaluation. MolTrustBench provides a reproducible scaffold for replacing implicit benchmark assumptions with explicit public-exposure lower bounds, assay-provenance diagnostics, and exposure-adjusted reporting.

\section{Data availability}

Processed evidence tables, exposure annotations, generated benchmark splits, figure source data, supplement tables, trust-card JSON examples, and analysis manifests underlying this article will be available from a DOI-issuing repository at \textbf{[repository DOI/accession to be inserted]}. The corresponding local evidence archive pending repository deposit is \artifact{moltrustbench_data_freeze_20260529_161700.tgz}; its SHA256 checksum is recorded in the deposit manifest \artifact{moltrustbench_data_freeze_20260529_161700_manifest.csv} and in the supplement data-availability file. The supplement package includes compact source data, slice-validity tables, standardization audit tables, assay-provenance flag definitions, machine-readable trust-card schemas, and a figure source-data map. The archive excludes raw ChEMBL database dumps, local secrets, transient caches, and third-party model checkpoints. ChEMBL source releases are available from the official ChEMBL downloads page \cite{chemblDownloads2025}; benchmark datasets and public checkpoints should be cited through their original sources.

\section{Code availability}

MolTrustBench code is maintained in the project repository. The current version uses the release-index, exposure-audit, assay-provenance, model-family integration, label-shuffle negative-control, supplement-packaging, and final-figure scripts under \artifact{scripts/}, with reusable package modules under \artifact{src/moltrustbench/}. \textbf{SUBMISSION\_METADATA\_REQUIRED:} add the GitHub URL, commit SHA, and tagged release corresponding to the deposited evidence archive.

\section{Acknowledgments}

\textbf{SUBMISSION\_METADATA\_REQUIRED:} add acknowledgments after author and compute-resource details are finalized.

\section{Funding}

\textbf{SUBMISSION\_METADATA\_REQUIRED:} add funding information or state that no specific funding was received.

\section{Conflict of interest}

\textbf{SUBMISSION\_METADATA\_REQUIRED:} verify the competing-interest statement before submission.

\bibliographystyle{oup-plain}
\bibliography{references}

\end{document}
